\par\medskip
\hypertarget{native-075-article057}{}\pdfbookmark[2]{57.}{article-057}
\begin{center}57.\end{center}
The distance of a celestial body from any plane passing through the Sun can thus be reduced to the form $kr\sin(v + K)$, where $v$ denotes the true anomaly, $k$ is the sine of the inclination of the orbit to this plane, and $K$ is the distance of perihelion from the ascending node of the orbit in that plane. Insofar as the orientation of the orbital plane, the position of the line of apsides in it, and the orientation of the reference plane for the distances can be regarded as constant, $k$ and $K$ will also be constant. This method will more often be used, however, in cases that rule out at least the third assumption, even if we neglect the perturbations that always affect the first two to some extent. This happens whenever the distances are referred to the
\SourceLeaf{076}equator or to a plane intersecting the equator at right angles at a given right ascension. The orientation of the equator changes owing to the precession of the equinoxes and also to nutation if we are considering the true position in place of the mean. In this case, therefore, $k$ and $K$ will also be subject to changes, which are of course slow.
These changes can be calculated by differential formulas that are easily derived. For brevity, it will suffice here to give the differential changes in $i'$, $\AscendingNode'$,
$\Delta$ insofar as they depend on changes in $\AscendingNode - n$
and $\epsilon$.
\[
\mathrm{d}i' = \sin\epsilon \sin\AscendingNode'\, \mathrm{d}(\AscendingNode - n) - \cos\AscendingNode'\, \mathrm{d}\epsilon
\]
\[
\mathrm{d}\AscendingNode' = \frac{\sin i \cos\Delta}{\sin i'}\, \mathrm{d}(\AscendingNode - n) + \frac{\sin\AscendingNode'}{\operatorname{tang} i'}\, \mathrm{d}\epsilon
\]
\[
\mathrm{d}\Delta = \frac{\sin\epsilon \cos\AscendingNode'}{\sin i'}\, \mathrm{d}(\AscendingNode - n) + \frac{\sin\AscendingNode'}{\sin i'}\, \mathrm{d}\epsilon.
\]
Moreover, when the sole aim is to calculate several positions of a celestial body with respect to such changing planes over a moderate interval of time, such as one year, it will usually be most convenient to calculate the quantities $a$, $A$, $b$, $B$, $c$, $C$ directly at two epochs bracketing those positions. Their changes for each of the specified times can then be found from these values by simple interpolation.
